% TOP_PROBLEM: {"id": 3260, "title": "Oriented trees in n-chromatic digraphs", "category_id": 3, "category": {"id": 3, "name": "graph_theory", "display_name": "Graph Theory", "description": "Problems involving graphs, networks, and their properties.", "slug": "graph-theory", "order_index": 3, "created_at": "2026-07-31T15:26:25.671Z"}, "status": "open", "problem_number": "OPG-46167", "set_id": 12}
% FIRST_POSTED: 2026-10-01
% ENTRY_KIND: statement_only
% UnsolvedMath ID 3260; source catalogue code OPG-46167
% !TeX program = lualatex
% Definition and sources only; this file is not an LLM solution attempt.
\documentclass[11pt,a4paper]{article}
\usepackage[margin=25mm]{geometry}
\usepackage{fontspec}
\setmainfont{lmroman10-regular.otf}[BoldFont=lmroman10-bold.otf,
 ItalicFont=lmroman10-italic.otf,BoldItalicFont=lmroman10-bolditalic.otf]
\usepackage{amsmath,amssymb,mathtools}
\usepackage{unicode-math}
\setmathfont{latinmodern-math.otf}
\AtBeginDocument{\let\setminus\smallsetminus}
\usepackage{xurl,hyperref}
\hypersetup{colorlinks=true,urlcolor=blue}
\setcounter{secnumdepth}{0}
\setlength{\parindent}{0pt}
\setlength{\parskip}{5pt}
\setlength{\emergencystretch}{3em}
\begin{document}
\section*{Oriented trees in n-chromatic digraphs}\label{3260}
{\small UnsolvedMath ID 3260 · OPG-46167 · Graph Theory · OpenGarden}
\subsection{Definitions and mathematical statement}
Let $D=(V,A)$ be a finite loopless simple digraph, allowing opposite arcs
but not repeated arcs in one direction. Its underlying undirected graph
$U(D)$ has an edge $uv$ whenever at least one of $uv,vu$ is an arc.
The chromatic number used here is $\chi(D)=\chi(U(D))$: adjacent vertices
must have different colours, regardless of arc direction. It is not
the dichromatic number defined by acyclic colour classes.

An oriented tree $T$ is obtained by orienting each edge of a finite
undirected tree. Its order is $|V(T)|$. A copy of $T$ as a subdigraph of
$D$ means an injective map $\varphi:V(T)\to V(D)$ such that
$xy\in A(T)$ implies $\varphi(x)\varphi(y)\in A(D)$. Other arcs among
the image vertices are allowed, so this is not an induced-copy condition.

The conjecture asks whether, for every integer $k\ge2$, every digraph
$D$ with
\[
                            \chi(D)\ge2k-2
\]
contains every oriented tree of order $k$ in this sense. The quantifier
``every'' includes all tree shapes and all choices of directions.

The proposed threshold is sharp: a regular tournament on $2k-3$ vertices
has ordinary chromatic number $2k-3$ but every outdegree is $k-2$, so it
cannot contain the oriented star with one vertex pointing to $k-1$
leaves. No bound on indegrees, outdegrees, or the total order of $D$
is assumed beyond the chromatic hypothesis.
\subsection{Short English statement}
Does ordinary chromatic number at least 2k−2 force a digraph to contain every orientation of every k-vertex tree?
\subsection{Sources}
\begin{itemize}
\item[S1] ULAM.AI and the UnsolvedMath Contributors, supplied problems.json, record 3260 (OPG-46167), OpenGarden collection. Catalogue identity and source transcription; CC BY 4.0.
\url{https://huggingface.co/datasets/ulamai/UnsolvedMath}.
\item[S2] Open Problem Garden, Oriented trees in n-chromatic digraphs. Original problem and discussion; the mathematical formulation below is expanded from the supplied source record.
\url{http://www.openproblemgarden.org/op/oriented_trees_in_n_chromatic_digraphs}.
\end{itemize}
\end{document}
